% Part: history
% Chapter: biographies
% Section: gerhard-gentzen

\documentclass[../../../include/open-logic-section]{subfiles}

\begin{document}

\olfileid{his}{bio}{gen}

\olsection{গেরহার্ড গেৎসেন}

\olphoto{gentzen-gerhard}{গেরহার্ড গেৎসেন}

গঠনগত প্রমাণতত্ত্বের স্রষ্টা হিসেবে গেরহার্ড গেৎসেন
সবচেয়ে বেশি পরিচিত; বিশেষত স্বাভাবিক নিষ্পাদন ও
সিকোয়েন্ট কলনের !!{derivation} পদ্ধতি তিনিই নির্মাণ
করেন। জার্মানির গ্রাইফ্‌স্‌ভাল্ডে ১৯০৯ সালের ২৪
নভেম্বর তাঁর জন্ম হয়। প্রস্তুতিমূলক বিদ্যালয়ে ভর্তি
হওয়ার আগে তিন বছর তিনি বাড়িতেই পড়াশোনা করেন;
বিদ্যালয়ে তখন পড়াশোনার দিক থেকে অধিকাংশ সহপাঠীর
চেয়ে তিনি পিছিয়ে ছিলেন। তবু তিনি ছিলেন অসাধারণ
মেধাবী ছাত্র এবং গণিতে তাঁর বিশেষ দক্ষতা দেখা যায়।
তাঁর আগ্রহ ছিল নানা বিষয়ে; যেমন, মায়ের জন্য কবিতা
এবং বিদ্যালয়ের নাট্যমঞ্চের জন্য নাটকও লিখতেন।

গ্রাইফ্‌স্‌ভাল্ড বিশ্ববিদ্যালয়ে গেৎসেনের উচ্চশিক্ষা শুরু
হয়; পরে তিনি গ্যোটিঙেন, মিউনিখ ও বার্লিনে যান।
১৯৩৩ সালে হারমান ভাইলের অধীনে গ্যোটিঙেন
বিশ্ববিদ্যালয় থেকে তিনি ডক্টরেট লাভ করেন। (তাঁর
বেশির ভাগ গবেষণার তত্ত্বাবধান করেছিলেন পল
বের্নাইস; কিন্তু নাৎসিরা তাঁকে বিশ্ববিদ্যালয় থেকে
বরখাস্ত করে।) ১৯৩৪ সালে গেৎসেন ডাভিড
হিলবার্টের সহকারী হিসেবে কাজ শুরু করেন। সেই বছরই
তিনি \emph{Untersuchungen \"{u}ber das logische
Schlie\ss en I--II [যৌক্তিক নিষ্পাদন-সংক্রান্ত
অনুসন্ধান I--II]} প্রবন্ধদ্বয়ে সিকোয়েন্ট কলন ও
স্বাভাবিক নিষ্পাদনের !!{derivation} পদ্ধতি নির্মাণ
করেন। ১৯৩৬ সালে তিনি পেয়ানোর স্বতঃসিদ্ধগুলির
সঙ্গতি প্রমাণ করেন।

নাৎসিদের সঙ্গে গেৎসেনের সম্পর্ক জটিল। তাঁর পরামর্শদাতা
বের্নাইসকে যে সময়ে জার্মানি ছাড়তে বাধ্য করা হয়,
সেই সময়েই গেৎসেন নাৎসি আধাসামরিক সংগঠন
এসএ-র বিশ্ববিদ্যালয় শাখায় যোগ দেন। অনেক জার্মানের
মতো তিনিও নাৎসি দলের সদস্য ছিলেন। যুদ্ধের সময়
তিনি বিমান-গোয়েন্দা ইউনিটে টেলিযোগাযোগ কর্মকর্তা
হিসেবে কাজ করেন। তবে ১৯৪২ সালে স্নায়বিক
বিপর্যয়ের কারণে তাঁকে দায়িত্ব থেকে অব্যাহতি দেওয়া
হয়। নাৎসি দলের প্রতি তাঁর আনুগত্য ছিল কি না, নাকি
প্রাতিষ্ঠানিক সাফল্য নিশ্চিত করতেই তিনি দলে
যোগ দিয়েছিলেন, তা স্পষ্ট নয়।

১৯৪৩ সালে গেৎসেন প্রাগের জার্মান বিশ্ববিদ্যালয়ের
গণিত ইনস্টিটিউটে একটি পদের প্রস্তাব পান এবং তা
গ্রহণ করেন। কিন্তু ১৯৪৫ সালে প্রাগের নাগরিকেরা
জার্মান দখলদারির বিরুদ্ধে বিদ্রোহ করেন। সোভিয়েত
বাহিনী শহরে এসে বিশ্ববিদ্যালয়ের সব অধ্যাপককে
গ্রেপ্তার করে। নাৎসি সংগঠনগুলিতে সদস্যপদের কারণে
গেৎসেনকে একটি বাধ্যতামূলক শ্রমশিবিরে নিয়ে যাওয়া
হয়। ১৯৪৫ সালের ৪ আগস্ট ৩৫ বছর বয়সে নিজের
কারাকক্ষে অপুষ্টিতে তাঁর মৃত্যু হয়।

\begin{reading}
গেৎসেনের পূর্ণাঙ্গ জীবনী জানতে \citet{Menzler-Trott2007}
দেখো। নাৎসি শাসনাধীন গণিতবিদদের নিয়ে
\citet{Segal2014} একটি আকর্ষণীয় বই; তাতে গেৎসেনের
জীবন নিয়েও সংক্ষিপ্ত আলোচনা আছে। যৌক্তিক
নিষ্পাদন নিয়ে গেৎসেনের প্রবন্ধগুলি মূল জার্মান ভাষায়
পাওয়া যায় \citep{Gentzen1935a,Gentzen1935b}।
গেৎসেনের প্রবন্ধগুলির ইংরেজি অনুবাদ এক খণ্ডে
সংকলন করেছেন \citet{Gentzen1969}; তাতে একটি
সংক্ষিপ্ত জীবনীও রয়েছে।
\end{reading}

\end{document}
